\section{Discussion}\label{sec:discuss}

\subsection{What grouping still buys}\label{sec:asymmetry}

Grouping was introduced so that a goodness-of-fit test for logistic regression would have a reference
distribution. It also bounds the influence of a record whose prediction is wrong, because the record's
residual is divided by the variance of its clean groupmates. The tests that left grouping behind to
gain power lost that bound: the record-level tests lose their level with one corrupted record in a
thousand (Section~\ref{sec:c1}). So do grouped tests whose extreme groups are small or that group on the
covariates, and so do tests along bounded directions once the model is refitted
(Section~\ref{sec:whatprotects}).
Among the pooled tests compared here, EDGE is the most powerful on link and tail departures, and the Hosmer--Lemeshow statistic
at ten groups holds its level further under reversed records (Section~\ref{sec:edgefr}).

On clean data EDGE is comparable to Stukel's score test and the cubic calibration test.
Its advantage is that it keeps its level when a few records are wrong, at a modest cost in
power. The protection
has a size, which depends on the error and the partition (Section~\ref{sec:edgefr}), and it does not
extend to wrong outcomes, which a calibration test should detect. A grouped test that rejects while
covariate errors are suspected is therefore reporting something other than those errors; the same
test rejecting while outcome errors are suspected may be reporting them, and the data should be checked before the model is rebuilt. Tolerating a few
wrong inputs is the right behaviour when the question is whether the model is right, at development
and validation. Once the model is deployed, a patient whose input is wrong receives a wrong
prediction whatever the model's calibration; finding that patient is the job of checks at data
entry, not of a calibration test.

\subsection{Recommended use}\label{sec:whentouse}

\begin{enumerate}
\item Use the cubic basis in its unit form at the default partition, $G=\max(10,\lceil n/25\rceil)$, and
      report the calibration intercept and slope beside it; \texttt{edge.gof(fit)} reports it beside
      ten groups.
\item If the data are known to carry entry errors in the predictors, as routinely collected records
      may, plan ten groups instead, \texttt{edge.gof(fit, G = 10)}: they tolerate about ten exaggerated
      or five reversed records in $1000$ and ten reversed records in $5000$, at a cost of $0.017$ in mean
      power and about half the power against misfit confined to the extreme risks. Decide this before
      seeing the results. Whichever setting decides, if only the default partition rejects, check the
      records in the extreme groups. If the suspected errors are reversals only, up to about ten in
      $1000$, the symmetric basis at the default partition holds further; it fails against many
      exaggerated records. If protection matters more than power, use the Hosmer--Lemeshow statistic at
      ten groups. Beyond a few tens of corrupted records, clean the data; influence diagnostics do not
      find them (Section~\ref{sec:whatprotects}).
\item When the departure feared is a symmetric change in both tails, use the symmetric basis; on clean
      data with high discrimination (an area under the curve near $0.9$) or a strongly skewed
      covariate, in its score form, \texttt{edge.gof(fit, basis = "sym", weights = "score")}.
\item When a frozen model, logistic or not, is validated on new patients, use the external mode at ten
      groups, in which Proposition~\ref{prop:bounded} holds exactly because nothing is refitted:
      \texttt{run.all.external(y, p)}, which reports it beside Cox's recalibration test and the
      calibration intercept and slope. In the design studied its average power was that of Cox's test
      and its false-alarm rate stayed below $0.10$ where that test's did not; Stukel's and the cubic
      tests on the frozen linear predictor broke in the cohort experiment of Section~\ref{sec:realcorrupt}.
\end{enumerate}

\subsection{When another test is better}\label{sec:whennot}

Table~\ref{tab:whennot} lists the settings in which the study found a better choice. In one of them
the remedy is a different form of EDGE; in the others it is a different test. For each,
we compared the best of EDGE's forms with the rival named.

\begin{table}[htbp]\centering
\caption{Settings in which another test, or another form of EDGE, is the better choice.
Power comparisons are on clean data; the false-alarm rates quoted in the first two rows, and the last
two rows, are under corrupted records.}
\label{tab:whennot}
\small
\begin{tabular}{>{\raggedright\arraybackslash}p{0.26\textwidth}>{\raggedright\arraybackslash}p{0.21\textwidth}>{\raggedright\arraybackslash}p{0.44\textwidth}}\toprule
Setting & Better choice & Evidence \\ \midrule
Scarce events ($12\%$), clean data & GiViTI belt & The belt leads EDGE by up to $0.152$;
but it reads $0.089$ false alarms at one corrupted record in $1000$ and $0.397$ at ten. \\
A single smooth bend, clean data & Cubic calibration test or GiViTI belt & The cubic test leads in $15$
of $21$ asymmetric-link scenarios (mean $0.741$, belt $0.744$, EDGE $0.669$; largest gap $0.208$); but it
reads $0.127$
and $0.482$ false alarms at one and ten corrupted records in $1000$, where EDGE reads $0.056$ and
$0.097$. \\
A curve that crosses the diagonal twice & BAGofT or le Cessie's test & BAGofT rejects $155$ of $200$
replicates, EDGE $18$. \\
Rough or oscillating misfit & BAGofT & $1.000$ and $0.805$ on the fastest oscillation and the sawtooth,
against $0.438$ and $0.093$ for the best form of EDGE, the Stukel basis. \\
An omitted interaction & Le Cessie's test & Leads in $11$ of $15$ binary and $8$ of $12$ continuous
interaction scenarios. \\
Symmetric tails, or a strongly skewed covariate & Symmetric basis, score form & Stukel's one-parameter
score leads the default basis by $0.062$ on symmetric tails and the unit symmetric form by $0.252$ under
a skewed covariate; the score form ties it on both ($-0.002$ on average over $32$ scenarios, $-0.007$
to $+0.012$ over the four skewed ones, where its level is $0.049$ to $0.057$). \\
Many reversed records, protection first & Hosmer--Lemeshow at ten groups & $0.099$ against $0.182$ for
EDGE at ten groups with ten reversed records in $1000$, and $0.126$ against $0.249$ with
twenty-five in $5000$; its power is lower. \\
Tens of corrupted records & Clean the data & At fifty in $5000$ the default reads $0.740$ and ten groups
$0.107$; at a hundred no setting holds, and a robust fit does not help. \\
\bottomrule
\end{tabular}
\end{table}

\subsection{Limitations and further work}\label{sec:limits}

The study is a simulation over six departure families, with models of one to three covariates, five
in the external-validation study and twelve in one design, and one real cohort fitted on its routinely
recorded predictors. The corrupted-record results use one family of covariate error, a value multiplied by four or by
minus four, with three other mechanisms reported beside it; errors correlated with the outcome were not
examined. Robust tests were compared only in the form of Stukel's test and at $n=1000$, and the weighted grouped
tests of \citet{hosmer2002goodness} only with their general-purpose weight. Proposition~\ref{prop:bounded} holds at
fixed coefficients; when the model is refitted, the influence of a corrupted record through the fit is
that of the maximum-likelihood estimator, which is measured here and not bounded. And the number of
corrupted records the test tolerates is read from the measured surface: we had expected the protection
at the default partition to hold over the whole range $0.1\%$ to $0.5\%$, and at the upper end it does
not.

In external mode the test already applies to any model that reports probabilities; in the
development data the estimation adjustment $\Omega$ is derived here for logistic regression only, and
its form for other models is left for further work. A quantity that predicts the whole partition surface of
Figure~\ref{fig:partition} is still missing. The difference between wrong predictions and wrong
outcomes suggests reading a grouped and an ungrouped test together, since their disagreement carries
information about which kind of error is present. A frozen model can be monitored as patients arrive (Section~\ref{sec:stream}); a sequential reference
that uses the dependence between looks, rather than spending the level over them, is left for further
work.
